%%%%%%%%%%%%%%%%%%%%%%%%%%%%%%%%%%%%%%%%%%%%%%%%%%%%%%%%%%%%%%%%%%%%%%
% Lớp 11 | Hàm số lượng giác và phương trình lượng giác | Giá trị lượng giác của một góc lượng giác
% Dạng 1D1?2-4: Các góc lượng giác có liên quan đặc biệt: bù nhau, phụ nhau, đối nhau,
% Số câu: 56  (sinh tự động bởi tach_cau.py; từ nay thầy sửa trực tiếp file này)
%%%%%%%%%%%%%%%%%%%%%%%%%%%%%%%%%%%%%%%%%%%%%%%%%%%%%%%%%%%%%%%%%%%%%%

%%%%%%%%%%%%%%%% TRẮC NGHIỆM 4 PHƯƠNG ÁN %%%%%%%%%%%%%%%%

% nguồn: Ok_Toan11_TracNghiem_Lan2
\begin{ex}%[1D1N2-4]%[ma:3c7461a7]
	Giá trị biểu thức $P=\left[\tan  \dfrac{17\pi}{4}+\tan \left(\dfrac{7\pi}{2}-x\right)\right]^2+\left[\cot \dfrac{13\pi}{4}+\cot \left(7\pi-x\right)\right]^2$ bằng
	\choice 
		{$\dfrac{1}{\sin ^2x}$} 
		{$\dfrac{1}{\cos  ^2x}$}
		{\True $\dfrac{2}{\sin ^2x}$} 
		 {$\dfrac{2}{\cos ^2x}$}
	\loigiai{ 
		Ta có 
		\begin{itemize}
			\item  $\tan \dfrac{17\pi}{4}=\tan \left(\dfrac{\pi}{4}+4\pi \right)=\tan \dfrac{\pi}{4}=1$.
			 \item $\tan \left(\dfrac{7\pi}{2}-x\right)=\cot x$.
			\item  $\cot \dfrac{13\pi}{4}=\cot \left(\dfrac{\pi}{4}+3\pi \right)=\cot \dfrac{\pi}{4}=1$.
			\item $\cot \left(7\pi-x\right)=-\cot  x$.
		\end{itemize}
		Suy ra $P=\left(1+\cot x\right)^2+\left(1-\cot  x\right)^2=2+2\cot ^2x=\dfrac{2}{\sin ^2x}$.
	} 
\end{ex}

% nguồn: Ok_Toan11_TracNghiem_Lan2
\begin{ex}%[1D1N2-4]%[ma:ce5ab1a2]
	Rút gọn biểu thức $S=\cos \left(\dfrac{\pi}{2}-x\right)\sin \left(\pi-x\right)-\sin \left(\dfrac{\pi}{2}-x\right)\cos \left(\pi-x\right)$ ta  được
	\choice 
		 {$S=0$}
		{$S=\sin  ^2x-\cos ^2x$}
		{$S=2\sin  x\cos x$}
		{\True  $S=1$}
	 \loigiai{
		 Ta có
		\begin{align*}
			S&=\cos \left(\dfrac{\pi}{2}-x\right)\cdot \sin  \left(\pi-x\right)-\sin \left(\dfrac{\pi}{2}-x\right)\cdot \cos \left(\pi-x\right)\\
			 &=\sin x\cdot \sin x-\cos x\cdot \left(-\cos x\right)=\sin ^2x+\cos ^2x=1.
		\end{align*}
	 }
\end{ex}

% nguồn: Ok_Toan11_TracNghiem_Lan2
\begin{ex}%[1D1N2-4]%[ma:e11e6f10]
	Cho $P=\sin \left(\pi+\alpha \right)\cdot \cos \left(\pi-\alpha \right)$ và $Q=\sin \left(\dfrac{\pi}{2}-\alpha \right)\cdot \cos  \left(\dfrac{\pi}{2}+\alpha \right)$. Mệnh đề nào dưới đây là đúng?
	 \choice
		{\True  $P+Q=0$}
		 {$P+Q=-1$}
		{$P+Q=1$} 
		{$P+Q=2$} 
	\loigiai{ 
		Ta có 
		\begin{itemize}
			\item $P=\sin \left(\pi+\alpha \right)\cdot \cos  \left(\pi-\alpha \right)=-\sin \alpha\cdot \left(-\cos \alpha \right)=\sin \alpha\cdot \cos\alpha$.
			\item $Q=\sin \left(\dfrac{\pi}{2}-\alpha \right)\cdot \cos \left(\dfrac{\pi}{2}+\alpha \right)=\cos \alpha\cdot \left(-\sin  \alpha \right)=-\sin \alpha\cdot \cos \alpha$.
		\end{itemize}
		Khi đó $P+Q=\sin \alpha\cdot \cos \alpha-\sin \alpha\cdot \cos  \alpha=0$.
	} 
\end{ex}

% nguồn: Ok_Toan11_TracNghiem_Lan2
\begin{ex}%[1D1N2-4]%[ma:36417293]
	Đơn giản biểu thức $  A=\cos\left(\alpha-\dfrac{\pi}{2}\right)$, ta được
	\choice 
		{$A=\cos\alpha $} 
		{\True $A=\sin\alpha  $}
		{$A=\cos\alpha $} 
		{$A=-\sin\alpha $} 
	\loigiai{ 
	Ta  có \[A=\cos\left(\alpha-\dfrac{\pi}{2}\right)=\cos\left(\dfrac{\pi}{2}-\alpha\right)=\sin\alpha.\]}
\end{ex}

% nguồn: Ok_Toan11_TracNghiem_Lan2
\begin{ex}%[1D1N2-4]%[ma:f5bdfdaa]
	Giá  trị $\cot\dfrac{89\pi}{6}$ bằng
	 \choice
		{$\sqrt{3}$} 
		{\True  $-\sqrt{3}$}
		{$\dfrac{\sqrt{3}}{3}$} 
		{$-\dfrac{\sqrt{3}}{3}$} 
	 \loigiai{
	Ta  có \[\cot\dfrac{89\pi}{6}=\cot\left(14\pi+\dfrac{5\pi}{6}\right)=\cot\dfrac{5\pi}{6}=-\sqrt{3}.\]}
\end{ex}

% nguồn: Ok_Toan11_TracNghiem_Lan2
\begin{ex}%[1D1N2-4]%[ma:dfa5ecc6]
	Với mọi $\alpha $ thì $\sin\left(\dfrac{3\pi}{2}+\alpha\right)$ bằng 
	 \choice
		{$-\sin\alpha  $}
		{\True  $-\cos\alpha $}
		 {$\cos\alpha $}
		 {$\sin\alpha $}
	\loigiai{ 
	 Ta có \[\sin\left(\dfrac{3\pi}{2}+\alpha\right)=\sin\left(2\pi+\alpha-\dfrac{\pi}{2}\right)=\sin\left(\alpha-\dfrac{\pi}{2}\right)=-\sin\left(\dfrac{\pi}{2}-\alpha\right)=-\cos\alpha .\]}
\end{ex}

% nguồn: Ok_Toan11_TracNghiem_Lan2
\begin{ex}%[1D1N2-4]%[ma:b424ae21]
	Tìm khẳng định \textbf{sai}?  (với điều kiện các hệ thức đã xác định)
	 \choice
		 {$\tan (\pi +\alpha)=\tan \alpha $}
		{\True $\cos  \left(\dfrac{\pi}{2} +\alpha \right)=\sin \alpha $}
		{$\cot  (-\alpha)=-\cot \alpha $}
		{$\sin  (\pi -\alpha)=\sin \alpha $}
	 \loigiai{
	Khẳng định $\cos \left(\dfrac{\pi}{2} +\alpha \right)=\sin \alpha  $ là sai vì $\cos \left(\dfrac{\pi}{2} +\alpha \right)=-\sin \alpha $.}
\end{ex}

% nguồn: Ok_Toan11_TracNghiem_Lan2
\begin{ex}%[1D1N2-4]%[ma:6e83d96a]
	Đơn giản biểu thức $A=\cos  \left(\alpha -\dfrac{\pi}{2} \right)$, ta được
	\choice 
		{$\cos \alpha  $}
		 {\True $\sin \alpha $}
		{$- \cos  \alpha $}
		{$-\sin \alpha  $}
	 \loigiai{
	 Ta có $A=\cos \left(\alpha -\dfrac{\pi}{2} \right)$$=\cos \left(\dfrac{\pi}{2} -\alpha \right)$$=\sin \alpha $.}
\end{ex}

% nguồn: Ok_Toan11_TracNghiem_Lan2
\begin{ex}%[1D1N2-4]%[ma:21426de1]
	 Giá trị của biểu thức $S=3-\sin ^290^{\circ}+2\cos ^260^{\circ}-3\tan ^245^{\circ}$ bằng
	 \choice
		{$\dfrac{1}{2}$} 
		{\True  $-\dfrac{1}{2}$}
		{$1$} 
		{$3$} 
	\loigiai{ 
		Ta có $S=3-\sin ^290^{\circ}+2\cos ^260^{\circ}-3\tan ^245^{\circ}$  $=3-1^2+2\cdot \left(\dfrac{1}{2} \right)^2-3\cdot 1^2$ $=-\dfrac{1}{2}$.
	} 
\end{ex}

% nguồn: Ok_TracNghiemLop11_Lan1
\begin{ex}%[1D1N2-4]%[ma:a3352b4f]
	Tính giá trị của  $\cos \left[\dfrac{\pi }{3}+\left( 2k+1 \right)\pi \right]$ với $k$ là số nguyên.
	 \choice
		{$\cos \left[\dfrac{\pi  }{3}+\left( 2k+1 \right)\pi \right]=-\dfrac{\sqrt{3}}{2}$}
		{$\cos \left[\dfrac{\pi }{3}+\left( 2k+1 \right)\pi  \right]=\dfrac{1}{2}$}
		{\True $\cos \left[\dfrac{\pi  }{3}+\left( 2k+1 \right)\pi \right]=-\dfrac{1}{2}$}
		{$\cos \left[\dfrac{\pi }{3}+\left( 2k+1 \right)\pi  \right]=\dfrac{\sqrt{3}}{2}$}
	\loigiai{ 
	Ta có $\cos \left[\dfrac{\pi }{3}+\left( 2k+1 \right)\pi \right]=\cos \left(  \dfrac{\pi }{3}+\pi +k2\pi \right)=\cos \left( \dfrac{\pi }{3}+\pi \right)=-\cos \dfrac{\pi }{3}=-\dfrac{1}{2}.$}
\end{ex}

% nguồn: Ok_Toan11_TracNghiem_Lan2
\begin{ex}%[1D1H2-4]%[ma:29dcd3c2]
	Rút gọn biểu thức $A=\dfrac{\sin\left(-328^\circ\right)\cdot \sin{958^\circ}}{\cot 572^\circ  }-\dfrac{\cos\left(-508^\circ\right)\cdot \cos\left(-1022^\circ\right)}{\tan\left(-212^\circ\right)}$ ta được
	\choice 
		 {\True $A=-1$}
		 {$A= 1$}
		{$A= 0$} 
		 {$A= 2$}
	\loigiai{Ta có 
		\allowdisplaybreaks\begin{eqnarray*}
		A&=&\dfrac{\sin\left(-328^\circ\right)\cdot\sin{958^\circ}}{\cot 572^\circ }-\dfrac{\cos\left(-508^\circ\right)\cdot  \cos\left(-1022^\circ\right)}{\tan\left(-212^\circ\right)}\\
		&=&-\dfrac{\sin 32^\circ\cdot\sin 58^\circ}{\cot 32^\circ}-\dfrac{\cos 32^\circ\cdot\cos 58^\circ}{\tan  32^\circ}\\
		&=&-\dfrac{\sin 32^\circ\cdot\cos 32^\circ}{\cot  32^\circ}-\dfrac{\cos 32^\circ\cdot \sin 32^\circ }{\tan 32^\circ}\\
		 &=&-\sin^232^\circ-\cos^232^\circ\\
		&=&-1. 
	\end{eqnarray*}
} 
\end{ex}

% nguồn: Ok_Toan11_TracNghiem_Lan2
\begin{ex}%[1D1H2-4]%[ma:27276765]
	Biết $A,$ $B$, $C$ là các góc của tam giác $ABC$. Khi  đó khẳng định nào sau đây đúng?
	 \choice
		{$\sin C=-\sin \left(A+B\right)$} 
		{$\cos C=\cos  \left(A+B\right)$}
		{$\tan C=\tan \left(A+B\right)$} 
		{\True $\cot  C=-\cot \left(A+B\right)$}
	 \loigiai{
		Vì $A$, $B$, $C$ là các góc của tam giác $ABC$ nên $A+B+C=180^\circ \Rightarrow C=180^\circ  -\left(A+B\right)$.\\
		Do đó $C$ và $\left(A+B\right)$ là 2 góc bù  nhau.
		\[\Rightarrow \sin C=\sin \left(A+B\right).\] 
	 }
\end{ex}

% nguồn: Ok_Toan11_TracNghiem_Lan2
\begin{ex}%[1D1H2-4]%[ma:286470c3]
	Đơn giản biểu  thức $ D=\sin\left(\dfrac{5\pi}{2}-\alpha\right)+\cos\left(13\pi+\alpha\right)-3\sin\left(\alpha-5\pi\right)$.
	\choice 
		{$ 3\sin\alpha-2\cos\alpha  $}
		{\True $ 3\sin\alpha $} 
		 {$-3\sin\alpha $}
		{$ 2\cos\alpha+3\sin\alpha $} 
	\loigiai{ 
		Ta có 
		\allowdisplaybreaks\begin{eqnarray*}
		D&=&\sin\left(\dfrac{5\pi}{2}-\alpha\right)+\cos\left(13\pi+\alpha\right)-3\sin\left(\alpha-5\pi\right)\\ 
		&=&\sin\left(\dfrac{\pi}{2}-\alpha\right)+\cos\left(\pi+\alpha\right)+3\sin\left(\pi-\alpha\right)\\ 
		&=& \cos\alpha-\cos\alpha+3\sin\alpha\\ 
		&=& 3\sin\alpha. 
	\end{eqnarray*}
} 
\end{ex}

% nguồn: Ok_Toan11_TracNghiem_Lan2
\begin{ex}%[1D1H2-4]%[ma:ac2da463]
	Đơn giản biểu thức $ D=\sin\left(\dfrac{5\pi}{2}-a\right)+\cos\left(13\pi+a\right)-3\sin\left(a-5\pi\right)$ ta  được
	\choice 
		{$D= 2\cos  a+3\sin a$}
		{$D= 3\sin  a-2\cos a$}
		{\True $D=3\sin a$} 
		{ $D= 4\cos  a-\sin a$}
	\loigiai{ 
		 Ta có
		\allowdisplaybreaks\begin{eqnarray*}
		D&=&\sin\left(\dfrac{5\pi}{2}-a\right)+\cos\left(13\pi+a\right)-3\sin\left(a-5\pi\right)\\ 
		&=&  \sin\left(2\pi+\dfrac{\pi}{2}-a\right)+\cos\left(12\pi+\pi+a\right)-3\sin\left(a+\pi-6\pi\right)\\
		&=&\sin\left(\dfrac{\pi}{2}-a\right)+\cos\left(\pi+a\right)-3\sin\left(a+\pi\right)\\ 
		 &=&\cos a-\cos a+3\sin a\\
		 &=&3\sin a.
	\end{eqnarray*}
 }
\end{ex}

% nguồn: Ok_Toan11_TracNghiem_Lan2
\begin{ex}%[1D1H2-4]%[ma:3cf9c5bc]
	Cho $\dfrac{\pi}{2} < \alpha < \pi$. Xét các mệnh đề  sau đây:
	\begin{enumEX}[I.]{3}
		\item $\cos \left(\alpha+\dfrac{\pi}{2} \right) < 0$. 
		\item $\sin \left(\alpha+\dfrac{\pi}{2} \right) < 0$. 
		\item $\cot \left(\alpha+\dfrac{\pi}{2} \right)  > 0$.
	\end{enumEX}
	Mệnh đề nào đúng? 
	\choice 
		{Chỉ I} 
		 {Chỉ I và II}
		{Chỉ II  và III}
		{\True Cả I,  II và III}
	 \loigiai{
		Ta có $\dfrac{\pi}{2} < \alpha < \pi \Rightarrow \pi < \left(\alpha+\dfrac{\pi}{2} \right) <  \dfrac{3\pi}{2}$ nên mệnh đề I, II và III.
	} 
\end{ex}

% nguồn: Ok_Toan11_TracNghiem_Lan2
\begin{ex}%[1D1H2-4]%[ma:4458d2da]
	Giá trị của biểu thức $\dfrac{\tan 225^\circ-\cot 81^\circ \cdot \cot 69^\circ}{\cot  261^\circ+\tan 201^\circ}$ bằng
	 \choice
		{$\dfrac{1}{\sqrt{3}}$} 
		 {$-\dfrac{1}{\sqrt{3}}$}
		{\True  $\sqrt{3}$}
		 {$-\sqrt{3}$}
	\loigiai{ 
		Ta  có
		\allowdisplaybreaks\begin{eqnarray*}
		\dfrac{\tan 225^\circ-\cot 81^\circ \cdot \cot 69^\circ}{\cot 261^\circ+\tan 201^\circ}&=&\dfrac{\tan\left(180^\circ+45^\circ\right)-\tan 9^\circ  \cdot \cot 69^\circ}{\cot\left(180^\circ+81^\circ\right)+\tan\left(180^\circ+21^\circ\right)}\\
		&=& \dfrac{1-\tan 9^\circ \cdot \tan 21^\circ}{\tan  9^\circ+\tan 21^\circ}\\
		 &=&\dfrac{1}{\tan\left(9^\circ+21^\circ\right)}\\
		 &=&\dfrac{1}{\tan 30^\circ}\\
		 &=& \sqrt{3}.
	\end{eqnarray*}
 }
\end{ex}

% nguồn: Ok_Toan11_TracNghiem_Lan2
\begin{ex}%[1D1H2-4]%[ma:186a959e]
	Nếu $\sin^2\alpha=\dfrac{1}{3}$ thì $ 1+\tan^2\alpha $ bằng 
	\choice 
		{$\dfrac{9}{8}$} 
		 {$4$}
		{\True $\dfrac{3}{2}$} 
		{$\dfrac{8}{9}$} 
	\loigiai{ 
	Ta có $\cos^2\alpha=1-\sin^2\alpha=\dfrac{2}{3}$ mà $ 1+\tan^2\alpha=\dfrac{1}{\cos^2\alpha}  \Rightarrow 1+\tan^2\alpha=\dfrac{3}{2}$.}
\end{ex}

% nguồn: Ok_Toan11_TracNghiem_Lan2
\begin{ex}%[1D1H2-4]%[ma:4d22c92c]
	Nếu $\cot(1{,}25)\cdot\tan \left(4\pi+1{,}25\right)-\sin \left(x+\dfrac{\pi}{2} \right)\cdot\cos \left(6\pi-x\right)=0$ thì $\tan x$  bằng
	\choice 
		 {$1$}
		 {$-1$}
		 {\True $0$}
		{$\dfrac{1}{2}$} 
	 \loigiai{
		Ta  có
		 \allowdisplaybreaks
		\begin{eqnarray*}
			&&  \cot (1{,}25)\cdot\tan \left(4\pi+1{,}25\right)-\sin \left(x+\dfrac{\pi}{2} \right)\cdot\cos \left(6\pi-x\right)=0 \\
			&\Leftrightarrow& \cot (1{,}25)\cdot\tan(1{,}25)-\cos x\cdot \cos  x=0 \\
			 &\Leftrightarrow& 1-\cos^2x=0 \\
			 &\Leftrightarrow& \sin^2x=0 \\
			&\Leftrightarrow& \sin x=0 \\ 
			 &\Leftrightarrow& \tan x=0.
		\end{eqnarray*}
	 }
\end{ex}

% nguồn: Ok_Toan11_TracNghiem_Lan2
\begin{ex}%[1D1H2-4]%[ma:2e867cc3]
	Với $\alpha$ là số  thực bất kỳ, rút gọn biểu thức $A=\cos \left(\alpha-\dfrac{\pi}{2}\right)+\sin (\alpha-\pi)$.
	\choice 
		 {$A=2 \sin \alpha$}
		{$A=2 \cos  \alpha$}
		{$A=1$} 
		 {\True $A=0$}
	 \loigiai{
		Ta có  $\cos \left(\alpha-\dfrac{\pi}{2}\right)=\cos \left(\dfrac{\pi}{2}-\alpha\right)=\sin \alpha$; $\sin (\alpha-\pi)=-\sin \alpha$.\\
	 Do đó $A=0$.}
\end{ex}

% nguồn: Ok_Toan11_TracNghiem_Lan2
\begin{ex}%[1D1H2-4]%[ma:4bb064e8]
	Đơn  giản biểu thức $ C=\cos\left(\dfrac{3\pi}{2}-a\right)-\sin\left(\dfrac{3\pi}{2}-a\right)+\cos\left(a-\dfrac{7\pi}{2}\right)-\sin\left(a-\dfrac{7\pi}{2}\right)$ ta được
	\choice 
		{$C=  2\sin a$}
		{\True $C=-2\sin  a$}
		 {$C= 2\cos a$}
		 {$C=-2\cos a$}
	 \loigiai{
		\allowdisplaybreaks\begin{eqnarray*}
		 C&=&\cos\left(\dfrac{3\pi}{2}-a\right)-\sin\left(\dfrac{3\pi}{2}-a\right)+\cos\left(a-\dfrac{7\pi}{2}\right)-\sin\left(a-\dfrac{7\pi}{2}\right)\\
		&=&\cos\left(2\pi-\dfrac{\pi}{2}-a\right)-\sin\left(2\pi-\dfrac{\pi}{2}-a\right)+\cos\left(a-4\pi+\dfrac{\pi}{2}\right)-\sin\left(a-4\pi+\dfrac{\pi}{2}\right)\\ 
		&=&\cos\left(\dfrac{\pi}{2}+a\right)-\sin\left(-\dfrac{\pi}{2}-a\right)+\cos\left(a+\dfrac{\pi}{2}\right)-\sin\left(a+\dfrac{\pi}{2}\right)\\ 
		&=&-\sin a+\cos a-\sin a-\cos a\\ 
		&=&-2\sin a. 
	\end{eqnarray*}
} 
\end{ex}

% nguồn: Ok_Toan11_TracNghiem_Lan2
\begin{ex}%[1D1H2-4]%[ma:a5d425d3]
	 Tính $ F=\sin^2\dfrac{\pi}{6}+\sin^2\dfrac{2\pi}{6}+\ldots +\sin^2\dfrac{5\pi}{6}+\sin^2\pi$.
	 \choice
		{\True $3$} 
		{$2$} 
		{$1$} 
		 {$4$}
	\loigiai{ 
		Ta có 
		\allowdisplaybreaks\begin{eqnarray*}
		F&=&\sin^2\dfrac{\pi}{6}+\sin^2\dfrac{2\pi}{6}+\ldots+\sin^2\dfrac{5\pi}{6}+\sin^2\pi\\ 
		&=&  \sin^2\dfrac{\pi}{6}+\sin^2\dfrac{\pi}{3}+\sin^2\dfrac{\pi}{2}+\sin^2\dfrac{2\pi}{3}+\sin^2\dfrac{5\pi}{6}+\sin^2\pi.\\
		&=&2\left(\sin^2\dfrac{\pi}{3}+\cos^2\dfrac{\pi}{3}\right)+1+0\\ 
		&=&3. 
	\end{eqnarray*}
 }
\end{ex}

% nguồn: Ok_Toan11_TracNghiem_Lan2
\begin{ex}%[1D1H2-4]%[ma:d11394de]
	Tính $S=\sin ^25^{\circ}+\sin ^210^{\circ}+\sin  ^215^{\circ}+\cdots\sin ^28^0+\sin ^285^{\circ}$.
	 \choice
		{$\dfrac{19}{2}$} 
		{$8$} 
		 {\True $\dfrac{17}{2}$}
		{$9$} 
	 \loigiai{
		\begin{align*}
			S&=\sin ^25^{\circ}+\sin ^210^{\circ}+\sin ^215^{\circ}+\cdots\sin ^280^{\circ}+\sin  ^285^{\circ}\\
			&=\left(\sin ^25^{\circ}+\sin  ^285^{\circ}\right)+\left(\sin ^210^{\circ}+\sin ^280^{\circ}\right)+\cdots\left(\sin ^215^{\circ}+\sin ^275^{\circ}\right)+\sin ^245^{\circ}\\
			&=8+\dfrac{1}{2}=\dfrac{17}{2}. 
		\end{align*}
	} 
\end{ex}

% nguồn: Ok_Toan11_TracNghiem_Lan2
\begin{ex}%[1D1H2-4]%[ma:616644ee]
	Cho $\dfrac{\pi}{2} < \alpha < \pi$.  Xét các mệnh đề sau:
	\begin{enumEX}[I.]{3}
		\item $\cos \left(\dfrac{\pi}{2}-\alpha  \right) > 0$.
		\item $\sin  \left(\dfrac{\pi}{2}-\alpha \right) > 0$.
		\item  $\tan \left(\dfrac{\pi}{2}-\alpha \right) > 0$.
	\end{enumEX}
	 Mệnh đề nào \textbf{sai}?
	\choice 
		 {Chỉ I}
		{Chỉ  II}
		 {\True Chỉ II và III}
		{Cả I, II và  III}
	 \loigiai{
		Ta có $\dfrac{\pi}{2} < \alpha < \pi \Rightarrow-\dfrac{\pi}{2}  < \alpha < 0$.
		Khi đó $\alpha$ thuộc cung phần tư thứ IV nên mệnh đề II,  III sai.
	 }
\end{ex}

% nguồn: Ok_Toan11_TracNghiem_Lan2
\begin{ex}%[1D1H2-4]%[ma:8ddcb394]
	Rút gọn biểu thức $ A=\dfrac{\sin\left(-234^\circ \right)-\cos216^\circ }{\sin 144^\circ -\cos126^\circ }\cdot \tan36^\circ  $, ta được
	\choice 
		{\True $  A=2$}
		{$ A=-2$} 
		 {$ A=1$}
		{$  A=-1$}
	 \loigiai{
		 Ta có
		\allowdisplaybreaks\begin{eqnarray*}
		A&=&\dfrac{\sin\left(-234^\circ \right)-\cos216^\circ }{\sin 144^\circ -\cos126^\circ }\cdot  \tan36^\circ \\
		&=&\dfrac{-\sin\left(180^\circ+54^\circ\right)-\cos\left(180^\circ +36^\circ \right)}{\sin\left(180^\circ  -36^\circ\right)-\cos\left(90^\circ+36^\circ\right)}\cdot \tan36^\circ\\
		&=&\dfrac{\sin 54^\circ+\cos36^\circ}{\sin 36^\circ+\sin  36^\circ}\cdot \tan 36^\circ\\
		&=&2\cot 36^\circ \cdot  \tan 36^\circ\\
		 &=&2.
	\end{eqnarray*}
 }
\end{ex}

% nguồn: Ok_Toan11_TracNghiem_Lan2
\begin{ex}%[1D1H2-4]%[ma:26f6618c]
	Giá trị của biểu thức $\tan 110^\circ\cdot \tan 340^\circ+\sin 160^\circ\cdot \cos 110^\circ+\sin 250^\circ\cdot \cos 340^\circ $  bằng
	 \choice
		 {\True $ 0$}
		{$ 1$} 
		{$-1$} 
		{$ 2$} 
	 \loigiai{
		 Ta có
		\allowdisplaybreaks\begin{eqnarray*}
		A&=&\tan 110^\circ\cdot \tan 340^\circ+\sin 160^\circ\cdot \cos  110^\circ+\sin 250^\circ\cdot \cos 340^\circ \\
		&=&\tan\left(90^\circ+20^\circ\right)\tan\left(360^\circ-20^\circ\right)+\sin\left(180^\circ-20^\circ\right)\cos\left(90^\circ+20^\circ\right)+\sin 250^\circ\cdot \cos 340^\circ \\ 
		&=&\cot 20^\circ\tan  20^\circ-\sin 20^\circ\sin 20^\circ+\sin\left(360^\circ-110^\circ\right)\cos\left(360^\circ-20^\circ\right)\\
		 &=&\cot 20^\circ\tan 20^\circ-\sin 20^\circ\sin 20^\circ-\sin 110^\circ\cos 20^\circ \\
		 &=&1-\sin^220^\circ-\sin\left(90^\circ+20^\circ\right)\cos 20^\circ\\
		&=&1-\sin^220^\circ-\cos^220^\circ  \\
		&=&1-\left(\sin^220^\circ+\cos^220^\circ\right)\\ 
		 &=&0.
	\end{eqnarray*}
} 
\end{ex}

% nguồn: Ok_Toan11_TracNghiem_Lan2
\begin{ex}%[1D1H2-4]%[ma:2ad47c04]
	Tính $ P=\cot 1^\circ \cdot \cot  2^\circ \cdot \cot 3^\circ \cdot \ldots \cdot \cot 89^\circ$.
	\choice 
		 {$P=0$}
		{\True  $P=1$}
		 {$P=3$}
		{$P=4$} 
	 \loigiai{
		 Ta có
		\allowdisplaybreaks\begin{eqnarray*}
		&&\cot 89^\circ=\tan 1^\circ \Rightarrow\cot 1^\circ\cdot  \cot 89^\circ=\cot 1^\circ\tan 1^\circ=1\\
		&&\cot 88^\circ=\tan 2^\circ \Rightarrow\cot 2^\circ\cdot \cot 88^\circ=\cot  2^\circ\tan 2^\circ=1\\
		 &&\ldots\\
		&&\cot 46^\circ=\tan 44^\circ \Rightarrow\cot  44^\circ\cdot \cot 46^\circ=\cot 44^\circ\tan 44^\circ=1.
	\end{eqnarray*}
Vậy $P=\cot 1^\circ \cdot \cot 2^\circ \cdot \cot 3^\circ \cdot \ldots \cdot \cot 89^\circ=\cot 45^\circ=1$.} 
\end{ex}

% nguồn: Ok_Toan11_TracNghiem_Lan2
\begin{ex}%[1D1H2-4]%[ma:84e9bd54]
	Giá trị  của biểu thức $A=\dfrac{1}{\tan 368^\circ}+\dfrac{2\sin 2550^\circ \cdot \cos\left(-188^\circ\right)}{2\cos 638^\circ +\cos 98^\circ}$ bằng
	 \choice
		 {$1$}
		 {$ 2$}
		{$-1$} 
		{\True  $0$}
	\loigiai{ 
		Ta  có
		\allowdisplaybreaks\begin{eqnarray*}
		A&=&\dfrac{1}{\tan 368^\circ}+\dfrac{2\sin 2550^\circ \cdot  \cos\left(-188^\circ\right)}{2\cos 638^\circ +\cos 98^\circ}\\
		& =& \dfrac{1}{\tan\left(8^\circ +360^\circ \right)}+\dfrac{2\sin\left(30^\circ+7\cdot  360^\circ\right)\cdot\cos\left(8^\circ+180^\circ\right)}{2\cos\left(-82^\circ+2\cdot 360^\circ\right)+\cos\left(90^\circ+8^\circ\right)}\\
		&=&\dfrac{1}{\tan 8^\circ}+\dfrac{-2\sin 30^\circ\cdot  \cos 8^\circ}{2\cos 82^\circ -\sin 8^\circ}\\
		&=&\dfrac{1}{\tan 8^\circ}-\dfrac{2\sin  30^\circ\cdot \cos 8^\circ}{2\cos\left(90^\circ-8^\circ\right)-\sin 8^\circ}\\
		&=&\dfrac{1}{\tan 8^\circ}-\dfrac{2\sin 30^\circ\cdot \cos 8^\circ}{2\sin 8^\circ-\sin  8^\circ}\\
		 &=&\cot 8^\circ-\dfrac{ \cos 8^\circ}{\sin 8^\circ}\\
		&=&\cot  8^\circ-\cot 8^\circ\\
		 &=&0.
	\end{eqnarray*}
 }
\end{ex}

% nguồn: Ok_Toan11_TracNghiem_Lan2
\begin{ex}%[1D1H2-4]%[ma:bb735335]
	Cho $\pi < \alpha < \dfrac{3\pi}{2}$. Xác định dấu của biểu thức $M=\sin \left(\dfrac{\pi}{2}-\alpha \right)\cdot\cot  \left(\pi+\alpha \right)$.
	 \choice
		{$M\ge 0$} 
		 {$M> 0$}
		{$M\le 0$} 
		{\True $M< 0$} 
	 \loigiai{
		Ta  có \\
		$\heva{&\pi < \alpha <  \dfrac{3\pi}{2} \\
		&{\pi <  \alpha < \dfrac{3\pi}{2}}}
		\Rightarrow \heva{&-\dfrac{3\pi}{2} <-\alpha <-\pi  \\ &2\pi < \pi+\alpha < \dfrac{5\pi}{2} }
		\Rightarrow \heva{&-\pi <  \dfrac{\pi}{2}-\alpha <-\dfrac{\pi}{2} \\ &\cot \left(\pi+\alpha \right) > 0}
		\Rightarrow \heva{&\sin \left(\dfrac{\pi}{2}-\alpha \right) < 0 \\  &\cot \left(\pi+\alpha \right) > 0}
		\Rightarrow M< 0.$ 
	 }
\end{ex}

% nguồn: Ok_Toan11_TracNghiem_Lan2
\begin{ex}%[1D1H2-4]%[ma:dc82af42]
	Với mọi góc $\alpha$, biểu thức  $\cos \alpha+\cos \left(\alpha+\dfrac{\pi}{5}\right)+\cos \left(\alpha+\dfrac{2\pi}{5}\right)+\cdots+\cos \left(\alpha+\dfrac{9\pi}{5}\right)$ nhận giá bằng
	 \choice
		{$10$} 
		 {$-10$}
		{$1$} 
		{\True  $0$}
	\loigiai{ 
		Trong tổng trên có $10$ số hạng,  chia thành $5$ cặp.\\
		 Vận dụng $\alpha+\dfrac{9\pi}{5}=\pi+\left(\alpha+\dfrac{4\pi}{5}\right) \Rightarrow \cos \left(\alpha+\dfrac{9\pi}{5}\right)=-\cos \left(\alpha+\dfrac{4\pi}{5}\right)$.\\
		 Tương tự
		\begin{itemize}
			\item $\cos \left(\alpha+\dfrac{8\pi}{5}\right)=-\cos \left(\alpha+\dfrac{3\pi}{5}\right)$. 
			 \item $\cos \left(\alpha+\dfrac{7\pi}{5}\right)=-\cos \left(\alpha+\dfrac{2\pi}{5}\right)$.
			\item $\cos \left(\alpha+\dfrac{6\pi}{5}\right)=-\cos \left(\alpha+\dfrac{\pi}{5}\right)$. 
			 \item $\cos \left(\alpha+\dfrac{5\pi}{5}\right)=-\cos \alpha$.
		\end{itemize}
		Do đó  $P=0$.
	} 
\end{ex}

% nguồn: Ok_Toan11_TracNghiem_Lan2
\begin{ex}%[1D1H2-4]%[ma:e44266a3]
	Cho $\dfrac{\pi}{2} < \alpha < \pi$. Xác định dấu  của biểu thức $M=\cos \left(-\dfrac{\pi}{2}+\alpha \right)\cdot\tan \left(\pi-\alpha \right)$.
	 \choice
		 {$M\ge 0$}
		{\True $M> 0$} 
		 {$M\le 0$}
		{$M< 0$} 
	\loigiai{ 
		 Ta có
		$\heva{&{\dfrac{\pi}{2} < \alpha < \pi } \\ &{\dfrac{\pi}{2} < \alpha  < \pi }}
		\Rightarrow \heva{&0<-\dfrac{\pi}{2}+\alpha < \dfrac{\pi}{2} \\  &0< \pi-\alpha < \dfrac{\pi}{2}}
		\Rightarrow \heva{&\cos\left(-\dfrac{\pi}{2}+\alpha \right)  > 0\\&\tan \left(\pi-\alpha \right) > 0}
		\Rightarrow M>  0$.
	} 
\end{ex}

% nguồn: Ok_Toan11_TracNghiem_Lan2
\begin{ex}%[1D1H2-4]%[ma:f5643751]
	Nếu $\cot(x+\pi)-\tan\left(x-\dfrac{\pi}{2}\right)=\sin^2(-1\,445^\circ)+\cos^2(1\,085^\circ)$ thì $\sin x$  bằng
	 \choice
		{$\pm\dfrac{1}{5}$} 
		 {$\pm\dfrac{2}{5}$}
		 {$\pm\dfrac{1}{\sqrt5}$}
		{\True $\pm\dfrac{2}{\sqrt5}$} 
	 \loigiai{
		Ta có 
		\allowdisplaybreaks 
		\begin{eqnarray*}
			&&\cot(x+\pi)-\tan\left(x-\dfrac{\pi}{2}\right)=\sin^2(-1\,445^\circ)+\cos^2(1\,085^\circ) \\ 
			 &\Leftrightarrow& \cot x - \tan\left[-\left(\dfrac{\pi}{2}-x\right)\right] = \sin^2(1\,445^\circ) +\cos^2(1\,085^\circ) \\
			&\Leftrightarrow& \cot x  +\tan\left(\dfrac{\pi}{2}-x\right) = \sin^2(1\,085^\circ + 360^\circ) + \cos^2(1\,085^\circ) \\
			&\Leftrightarrow& \cot x + \cot  x = \sin^2(1\,085^\circ) + \cos^2(1\,085^\circ) \\
			 &\Leftrightarrow& 2\cot x= 1 \\
			&\Leftrightarrow&  \cot x = \dfrac{1}{2}.
		\end{eqnarray*}
		Khi đó 
		\[\sin  x = \pm\dfrac{1}{\sqrt{1+\cot^2 x}} = \pm\dfrac{2}{\sqrt5}.\]
	} 
\end{ex}

% nguồn: Ok_TracNghiemLop11_Lan1
\begin{ex}%[1D1H2-4]%[ma:c36a49dd]
	Hãy chọn kết quả sai trong các kết quả sau  đây
	 \choice
		{$\cos\left(-\alpha\right)=\cos\alpha $} 
		{\True $\sin\left(\pi+\alpha\right)=\sin\alpha  $}
		 {$\tan\left(\pi-\alpha\right)=-\tan\alpha $}
		 {$\cot\left(\dfrac{\pi}{2}-\alpha\right)=\tan\alpha $}
	\loigiai{Ta có $\sin\left(\pi+\alpha\right)=-\sin\alpha $  nên $\sin\left(\pi+\alpha\right)=\sin\alpha $ {\bf sai}.}
\end{ex}

% nguồn: Ok_TracNghiemLop11_Lan1
\begin{ex}%[1D1H2-4]%[ma:fcbea0a7]
	 Tìm khẳng định {\bf sai} (với điều kiện các hệ thức đã xác định)?
	 \choice
		{$\tan\left(\pi+\alpha\right)=\tan\alpha $} 
		{\True  $\cos\left(\dfrac{\pi}{2}+\alpha\right)=\sin\alpha $}
		{$\cot\left(-\alpha\right)=-\cot\alpha $} 
		{$\sin\left(\pi-\alpha\right)=\sin\alpha  $}
	 \loigiai{
		$\cos\left(\dfrac{\pi}{2}+\alpha\right)=-\sin\alpha $ nên $\cos\left(\dfrac{\pi}{2}+\alpha\right)=\sin\alpha $ là khẳng định sai. 
	 }
\end{ex}

% nguồn: Ok_Toan11_TracNghiem_Lan2
\begin{ex}%[1D1V2-4]%[ma:4dc38413]
	Với mọi góc $\alpha$, biểu thức $\cos\alpha+\cos\left(\alpha+\dfrac{\pi}{5}\right)+\cos\left(\alpha+\dfrac{2\pi}{5}\right)+\ldots +\cos\left(\alpha+\dfrac{9\pi}{5}\right)$  nhận giá trị bằng
	 \choice
		{$  10$}
		{$-10$} 
		{$ 1$} 
		{\True $  0$}
	 \loigiai{
		Ta có 
		\allowdisplaybreaks\begin{eqnarray*}
		&&\cos\alpha=-\cos\left(\alpha+\dfrac{5\pi}{5}\right);\\ 
		&&\cos\left(\alpha+\dfrac{\pi}{5}\right)=-\cos\left(\alpha+\dfrac{6\pi}{5}\right);\\ 
		&&\cos\left(\alpha+\dfrac{2\pi}{5}\right)=-\cos\left(\alpha+\dfrac{7\pi}{5}\right);\\ 
		 &&\cos\left(\alpha+\dfrac{3\pi}{5}\right)=-\cos\left(\alpha+\dfrac{8\pi}{5}\right);\\
		 &&\cos\left(\alpha+\dfrac{4\pi}{5}\right)=-\cos\left(\alpha+\dfrac{9\pi}{5}\right).\\
	\end{eqnarray*}
Do  đó $\cos\alpha+\cos\left(\alpha+\dfrac{\pi}{5}\right)+\cos\left(\alpha+\dfrac{2\pi}{5}\right)+\ldots+\cos\left(\alpha+\dfrac{9\pi}{5}\right)=0$.}
\end{ex}

% nguồn: Ok_Toan11_TracNghiem_Lan2
\begin{ex}%[1D1V2-4]%[ma:c3c7639f]
	Cho $B=\dfrac{\cos ^2696^{\circ}+\tan \left(-260^{\circ}\right) \cdot \tan 530^{\circ}-\cos ^2156^{\circ}}{\tan ^2252^{\circ}+\cot  ^2342^{\circ}}$. Biểu thức thu gọn của $B$ là
	\choice 
		{$\dfrac{1}{2}  \tan ^224^0$}
		{$\dfrac{1}{2}  \cot ^224^0$}
		{\True $\dfrac{1}{2} \tan  ^218^0$}
		{$\dfrac{1}{2}  \cot ^218^0$}
	 \loigiai{
		Ta  có
		\[ 
		\begin{aligned}
			B&=\dfrac{\cos  ^2\left(720^{\circ}-24^{\circ}\right)-\tan \left(360^{\circ}-100^{\circ}\right) \cdot \tan \left(360^{\circ}+170^{\circ}\right)-\cos ^2\left(180^{\circ}-24^{\circ}\right)}{\tan ^2\left(360^{\circ}-108^{\circ}\right)+\cot ^2\left(360^{\circ}-18^{\circ}\right)}
			\\ 
			 &=\dfrac{\cos ^224^{\circ}+\tan \left(90^{\circ}+10^{\circ}\right) \cdot \tan \left(180^{\circ}-10^{\circ}\right)-\cos ^224^{\circ}}{\tan ^2\left(90^{\circ}+18^0\right)+\cot ^218^0}\\
			&=\dfrac{-\cot 10^{\circ} \cdot\left(-\tan  10^{\circ}\right)}{\cot ^218^{\circ}+\cot ^218^0}\\
			&=\dfrac{1}{2\cot ^218^0}\\ 
			&=\dfrac{1}{2} \tan  ^218^0.
		\end{aligned}
		 \]
	} 
\end{ex}

% nguồn: Ok_Toan11_TracNghiem_Lan2
\begin{ex}%[1D1V2-4]%[ma:2f7a449a]
	Biểu thức $A=\dfrac{\sin 515^{\circ} \cdot \cos \left(-475^{\circ}\right)+\cot 222^{\circ} \cdot \cot 408^{\circ}}{\cot 415^{\circ} \cdot \cot \left(-505^{\circ}\right)+\tan 197^{\circ}  \cdot \tan 73^{\circ}}$ có kết quả rút gọn bằng
	\choice 
		{$\dfrac{1}{2} \sin  ^225^{\circ}$}
		{$\dfrac{1}{2}  \cos ^255^{\circ}$}
		{\True  $\dfrac{1}{2} \cos ^225^{\circ}$}
		{$\dfrac{1}{2} \sin  ^265^{\circ}$}
	\loigiai{ 
		Ta có 
		\[ 
		\begin{aligned}
			A&=\dfrac{\sin 155^{\circ} \cdot \cos 115^{\circ}+\cot 42^{\circ} \cdot \cot 48^{\circ}}{\cot 55^{\circ} \cdot \cot \left(-145^{\circ}\right)+\tan 17^{\circ} \cdot \cot 17^{\circ}} \\ 
			&=\dfrac{\sin 25^{\circ} \cdot\left(-\sin 25^{\circ}\right)+\cot 42^{\circ} \cdot  \tan 42^{\circ}}{\cot 55^{\circ} \cdot \tan 55^{\circ}+1}\\
			&=\dfrac{-\sin  ^225^{\circ}+1}{2}\\
			&  =\dfrac{\cos ^225^{\circ}}{2}.
		\end{aligned}
		 \]
	 }
\end{ex}

% nguồn: Ok_Toan11_TracNghiem_Lan2
\begin{ex}%[1D1V2-4]%[ma:281d700d]
	Tính $  P=\sin \left(\alpha+\dfrac{\pi}{2}\right)+\cos (3\pi-2\alpha)+\cot (\pi-\alpha)$, biết $\sin \alpha=-\dfrac{1}{2}$ và $-\dfrac{\pi}{2} < \alpha < 0$.
	 \choice
		 {\True $\dfrac{3\sqrt{3}-1}{2}$}
		{$\dfrac{3\sqrt{3}-3}{2}$} 
		{$\dfrac{3\sqrt{3}+3}{2}$} 
		 {$\dfrac{3\sqrt{3}+1}{2}$}
	 \loigiai{
		\begin{itemize}
			 \item \textbf{Cách 1}
			\[ 
			\begin{aligned}
				P&=\sin  \left(\alpha+\dfrac{\pi}{2}\right)+\cos (3\pi-2\alpha)+\cot (\pi-\alpha)\\
				&=\cos (-\alpha)-\cos (-2\alpha)+\cot  (-\alpha)\\
				&=\cos \alpha-\cos 2\alpha-\cot  \alpha=\cos \alpha-\left(2\cos ^2\alpha-1\right)-\cot \alpha.\\
			\end{aligned}
			 \]
			Mặt khác $\cos ^2\alpha=1-\sin ^2\alpha=1-\left(-\dfrac{1}{2}\right)^2=\dfrac{3}{4}$ mà  $-\dfrac{\pi}{2} < \alpha < 0$ nên $\cos \alpha=\dfrac{\sqrt{3}}{2}$.\\
			Suy  ra $\cot \alpha=\dfrac{\cos \alpha}{\sin \alpha}=-\sqrt{3}$.\\
			Do đó $P=\cos \alpha-\left(2\cos  ^2\alpha-1\right)-\cot \alpha=\dfrac{\sqrt{3}}{2}-\left(2\cdot \dfrac{3}{4}-1\right)+\sqrt{3}=\dfrac{3\sqrt{3}-1}{2} $. \\
			\item \textbf{Cách 2}\\ 
			Vì $\sin \alpha=-\dfrac{1}{2}$ và $-\dfrac{\pi}{2} < \alpha < 0$  nên $\alpha=-\dfrac{\pi}{6}$.\\
			Thế vào $P$ ta được\\ 
			$P=\sin \left(-\dfrac{\pi}{6}+\dfrac{\pi}{2}\right)+\cos \left(3\pi+2\cdot  \dfrac{\pi}{6}\right)+\cot \left(\pi+\dfrac{\pi}{6}\right)=\sin \dfrac{\pi}{3}-\cos \dfrac{\pi}{3}-\cot \dfrac{\pi}{6}=\dfrac{3\sqrt{3}-1}{2}$.
		\end{itemize}
	 }
\end{ex}

% nguồn: Ok_Toan11_TracNghiem_Lan2
\begin{ex}%[1D1V2-4]%[ma:04640b75]
	Biểu  thức
	\[A=\cos (\alpha+26\pi)-2\sin (\alpha-7\pi)-\cos 1{,}5\pi-\cos \left(\alpha+\dfrac{2003\pi}{2}\right)+\cos (\alpha-1{,}5\pi) \cdot \cot (\alpha-8\pi)\] có kết quả thu gọn bằng 
	\choice 
		{$-\sin  \alpha$}
		{\True $\sin \alpha$} 
		 {$-\cos \alpha$}
		{$\cos  \alpha$}
	\loigiai{ 
		Ta có 
		\[ 
		\begin{aligned}
			 A&=\cos (\alpha+26\pi)-2\sin (\alpha-7\pi)-\cos (1{,}5\pi)-\cos \left(\alpha+2003\dfrac{\pi}{2}\right)+\cos (\alpha-1{,}5\pi) \cdot \cot (\alpha-8\pi)\\
			&=\cos \alpha-2\sin (\alpha-\pi)-\cos \left(\dfrac{\pi}{2}\right)-\cos \left(\alpha-\dfrac{\pi}{2}\right)+\cos \left(\alpha+\dfrac{\pi}{2}\right)  \cdot \cot \alpha\\
			&=\cos \alpha+2\sin \alpha-0-\sin \alpha-\sin \alpha \cdot \cot  \alpha\\
			&=\cos \alpha+\sin \alpha-\cos \alpha\\ 
			&=\sin  \alpha.
		\end{aligned}
		 \]
	} 
\end{ex}

% nguồn: Ok_Toan11_TracNghiem_Lan2
\begin{ex}%[1D1V2-4]%[ma:a86c72bb]
	Biểu thức \[A=\left[\tan (\pi-x)  \cdot \tan \left(\dfrac{3\pi}{2}+x\right) \cdot \dfrac{1}{\cos ^2\left(x-\dfrac{3\pi}{2}\right)}-\cos \left(\dfrac{3\pi}{2}+x\right) \cdot \dfrac{1}{\sin (\pi-x)}\right] \sin ^2(2\pi-x)\] có kết quả rút gọn bằng
	\choice 
		{$\sin ^2x$} 
		{\True  $\cos ^2x$}
		 {$\tan ^2x$}
		 {$\cot ^2x$}
	 \loigiai{
		Ta có 
		\[ 
		\begin{aligned}
			A &  ={\left[\tan (\pi-x) \cdot \tan \left(\dfrac{3 \pi}{2}+x\right) \dfrac{1}{\cos^2\left(x-\dfrac{3 \pi}{2}\right)}-\cos \left(\dfrac{3 \pi}{2}+x\right) \cdot \dfrac{1}{\sin (\pi-x)}\right] \sin^2(2 \pi-x)} \\
			&=\left[-\tan x \cdot \tan \left(\pi+\dfrac{\pi}{2}+x\right) \cdot \dfrac{1}{\cos^2\left(\pi+\dfrac{\pi}{2}-x\right)}-\cos \left(\pi+\dfrac{\pi}{2}+x\right) \dfrac{1}{\sin  x}\right] \sin^2 x \\
			&=\left[-\tan x \cdot(-\cot  x) \cdot \dfrac{1}{\sin^2 x}-\dfrac{\sin x}{\sin x} \cdot \sin^2 x\right. \\
			&=\left(\dfrac{1}{\sin^2 x}-1\right) \cdot \sin^2 x\\ 
			&=\cot^2 x \cdot \sin^2  x\\
			&=\cos^2  x.
		\end{aligned}
		\] 
	} 
\end{ex}

% nguồn: Ok_Toan11_TracNghiem_Lan2
\begin{ex}%[1D1B2-4]%[ma:a7c2e15e]
	 Cho $0< \alpha < \dfrac{\pi}{2}$. Khẳng định nào sau đây \textbf{đúng}?
	\choice 
		{$\sin \left(\alpha-\pi \right)\ge 0$} 
		{$\sin \left(\alpha-\pi \right)\le 0$} 
		{$\sin \left(\alpha-\pi  \right) > 0$}
		 {\True $\sin \left(\alpha-\pi \right) < 0$}
	 \loigiai{
		Ta có $0< \alpha < \dfrac{\pi}{2}  \Rightarrow -\pi < \alpha-\pi <-\dfrac{\pi}{2} $.\\
		Khi đó điểm cuối cung $\alpha-\pi$ thuộc góc phần tư thứ  III $ \Rightarrow \sin \left(\alpha-\pi \right) < 0$.
	} 
\end{ex}

% nguồn: Ok_Toan11_TracNghiem_Lan2
\begin{ex}%[1D1B2-4]%[ma:6dd1e8e3]
	Cho hai góc nhọn $\alpha$ và $\beta$ phụ nhau.  Hệ thức nào sau đây là \textbf{sai}?
	 \choice
		{\True $\sin \alpha=-\cos  \beta$}
		{$\cos \alpha=\sin \beta$} 
		{$\cos \beta=\sin  \alpha$}
		{$\cot  \alpha=\tan \beta$}
	 \loigiai{
		Do $\alpha$ và $\beta$  phụ nhau nên $\sin \alpha=-\cos \beta$.
	 }
\end{ex}


%%%%%%%%%%%%%%%% ĐÚNG/SAI %%%%%%%%%%%%%%%%

% nguồn: Ok_dungsai_lop11_Lan2
\begin{ex}%[1D1N2-4]%[ma:2473fd48]
	Cho $\alpha $  và $\beta $ là hai góc khác nhau và bù nhau.
	 \choiceTF
		 {$\cot \alpha =\cot \beta $}
		{\True $\sin \alpha  =\sin \beta $}
		{\True $\tan \alpha =-\tan  \beta $}
		{\True $\cos  \alpha =-\cos \beta $}
	 \loigiai{
		\begin{itemchoice}
			\itemch  \textbf{Sai.}\\
			Sửa cho đúng là $\cot \alpha =-\cot \beta $. 
			 \itemch \textbf{Đúng.}
			\itemch \textbf{Đúng.} 
			 \itemch \textbf{Đúng.}
		\end{itemchoice}
	} 
\end{ex}

% nguồn: Ok_dungsai_lop11_Lan1
\begin{ex}%[1D1H2-4]%[ma:59b2b5bf]
 Mỗi khẳng định  sau đúng hay sai?
 \choiceTF 
  {$\sin\left(\dfrac{2024\pi}{7}-x\right)=\sin\left(\dfrac{\pi}{7}-x\right)$} 
  {\True  $\tan\left(\dfrac{2023\pi}{2}+2\alpha\right)=-\cot2\alpha$}
   {\True $\cos\left(-\dfrac{2024\pi}{9}-x\right)=\cos\left(\dfrac{8\pi}{9}+x\right)$}
   {$\cot\left(\alpha+\dfrac{2024\pi}{13}\right)=-\cot\left(\alpha+\dfrac{4\pi}{13}\right)$}
 \loigiai{Ta có 
  \begin{itemchoice}
    \itemch $\sin\left(\dfrac{2024\pi}{3}-x\right)=\sin\left(288\pi+\pi+\dfrac{\pi}{7}-x\right)=\sin\left(\pi+\dfrac{\pi}{7}-x\right)=-\sin\left(\dfrac{\pi}{7}-x\right)$.
   \itemch $\tan\left(\dfrac{2023\pi}{2}+2\alpha\right)=\tan\left(1011\pi+\dfrac{\pi}{2}+2\alpha\right)=\tan\left(\dfrac{\pi}{2}+2\alpha\right)=-\cot2\alpha$. 
   \itemch $\cos\left(-\dfrac{2024\pi}{9}-x\right)=\cos\left(-224\pi-\dfrac{8\pi}{9}-x\right)=\cos\left(-\dfrac{8\pi}{9}-x\right)=\cos\left(\dfrac{8\pi}{9}+x\right)$. 
    \itemch $\cot\left(\alpha+\dfrac{2024\pi}{13}\right)=\cot\left(\dfrac{9\pi}{13}+\alpha\right)=\cot\left[\pi-\left(\dfrac{4\pi}{13}-\alpha\right)\right]=-\cot\left(\dfrac{4\pi}{13}-\alpha\right)$.
  \end{itemchoice}
  }
\end{ex}

% nguồn: Ok_dungsai_lop11_Lan1
\begin{ex}%[1D1H2-4]%[ma:3bcba0e4]
 Mỗi khẳng định sau đúng hay  sai?
  \choiceTF
   {$\sin\left(\dfrac{2024\pi}{7}-x\right)=\sin\left(\dfrac{\pi}{7}-x\right)$}
  {\True $\cos\left(-\dfrac{2024\pi}{9}-x\right)=\cos\left(\dfrac{8\pi}{9}+x\right)$} 
   {\True $\tan\left(\dfrac{2023\pi}{2}+2\alpha\right)=-\cot2\alpha$}
   {$\cot\left(\alpha+\dfrac{2024\pi}{13}\right)=-\cot\left(\alpha+\dfrac{4\pi}{13}\right)$}
 \loigiai{Ta có 
  \begin{itemchoice}
   \itemch  $\sin\left(\dfrac{2024\pi}{3}-x\right)=\sin\left(288\pi+\pi+\dfrac{\pi}{7}-x\right)=\sin\left(\pi+\dfrac{\pi}{7}-x\right)=-\sin\left(\dfrac{\pi}{7}-x\right)$.
    \itemch $\cos\left(-\dfrac{2024\pi}{9}-x\right)=\cos\left(-224\pi-\dfrac{8\pi}{9}-x\right)=\cos\left(-\dfrac{8\pi}{9}-x\right)=\cos\left(\dfrac{8\pi}{9}+x\right)$.
   \itemch  $\tan\left(\dfrac{2023\pi}{2}+2\alpha\right)=\tan\left(1011\pi+\dfrac{\pi}{2}+2\alpha\right)=\tan\left(\dfrac{\pi}{2}+2\alpha\right)=-\cot2\alpha$.
   \itemch $\cot\left(\alpha+\dfrac{2024\pi}{13}\right)=\cot\left(\dfrac{9\pi}{13}+\alpha\right)=\cot\left[\pi-\left(\dfrac{4\pi}{13}-\alpha\right)\right]=-\cot\left(\dfrac{4\pi}{13}-\alpha\right)$. 
  \end{itemchoice}
 } 
\end{ex}

% nguồn: Ok_dungsai_lop11_Lan1
\begin{ex}%[1D1H2-4]%[ma:5e0fb5b3]
 \immini{Cho đường tròn lượng giác gốc $A$ (tham khảo hình vẽ). Gọi $M$ là điểm biểu diễn góc  lượng giác có số đo $\alpha$. Mỗi khẳng định sau đúng hay sai?
 \choiceTF 
  {\True Nếu $M_1$ là điểm  đối xứng với điểm $M$ qua $Oy$ thì $M_1$ là điểm biểu diễn góc có số đo $(\pi-\alpha)$}
  {\True Nếu $M_2$ là điểm đối xứng với điểm $M$ qua $Ox$ thì $M_2$ là điểm  biểu diễn góc có số đo $-\alpha$}
  {\True Nếu $M_3$ là điểm đối xứng với điểm $M$ qua $O$ thì $M_3$ là điểm biểu diễn góc  có số đo $(\pi+\alpha)$}
  {\True Nếu $M_4$ là điểm đối xứng với điểm $M$ qua đường thẳng $y=x$ thì $M_4$ là điểm biểu diễn góc  có số đo $\left(\dfrac{\pi}{2}-\alpha\right)$ }
  } 
  { \begin{tikzpicture}[line join = round,line cap = round,>=stealth,font=\footnotesize,scale=0.77]
  \tikzset{label  style/.style={font=\footnotesize}}
  \path (0,0)  coordinate (O)
    (3,0) coordinate (A)
   (0,3)  coordinate (B)
   (0,-3) coordinate (B') 
    (-3,0) coordinate (A')
    (0:0)++(60:3) coordinate (M)
   ($(O)!(M)!(A')$) coordinate  (H)
   ($(O)!(M)!(B)$) coordinate (K); 
  \draw[->] (-4,0)  -- (4,0) node[above,blue]{$x$};
  \draw[->] (0,-4) -- (0.,4)  node[left,blue]{$y$};
  \draw[orange]  (O) circle (3cm);
  \draw[rotate=0,->,green!50!black]  (0.7,0) arc (0:60:0.7cm);
  \draw (.8,.5)  node[above,blue] {$\alpha$};
  \draw[green!50!black]  (M)--(O);
  \foreach \p/\r in  {A/-45,M/45,O/-45,A'/-135,B'/-45,B/45}
   \fill (\p) circle (1pt) node[shift={(\r:3mm)},blue]{$\p$};
 \end{tikzpicture}
  } 
  \loigiai{
 \immini{ 
  \begin{itemchoice}
   \itemch Nếu $M_1$ là điểm đối xứng với điểm $M$ qua $Oy$ thì $\widehat{M_1OA}=\pi-\alpha$ nên $M_1$ là điểm biểu  diễn góc có số đo $(\pi-\alpha)$.
   \itemch Nếu $M_2$ là điểm đối xứng với điểm $M$ qua $Ox$ thì $\widehat{M_2OA}=\alpha$ nên $M_2$ là  điểm biểu diễn góc có số đo $-\alpha$.
   \itemch Nếu $M_3$ là điểm đối xứng với điểm $M$ qua $O$ thì  $\widehat{M_3OA}=(\pi-\alpha)$ nên $M_3$ là điểm biểu diễn góc có số đo $(\pi+\alpha)$.
    \itemch
   Nếu $M_4$ là điểm đối xứng với điểm $M$ qua đường thẳng  $y=x$ thì $\widehat{M_4OA}=\left(\dfrac{\pi}{2}-\alpha\right)$ nên $M_4$ là điểm biểu diễn góc có số đo $\left(\dfrac{\pi}{2}-\alpha\right)$
  \end{itemchoice}}
  { 
  \begin{tikzpicture}[line join = round,line cap = round,>=stealth,font=\footnotesize,scale=0.77]
    \tikzset{label  style/.style={font=\footnotesize}}
    \path (0,0) coordinate (O) 
     (3,0) coordinate  (A)
     (0,3) coordinate  (B)
     (0,-3) coordinate  (B')
      (-3,0) coordinate (A')
     (0:0)++(60:3) coordinate (M) 
     (0:0)++(120:3) coordinate (M_1) 
     (0:0)++(-60:3)  coordinate (M_2)
     (0:0)++(-120:3) coordinate  (M_3)
     (0:0)++(30:3) coordinate (M_4) 
      ($(O)!(M)!(A')$) coordinate (H)
     ($(O)!(M)!(B)$)  coordinate (K);
    \draw[->] (-4,0) --  (4,0) node[above,blue]{$x$};
    \draw[->] (0,-4) -- (0.,4) node[left,blue]{$y$}; 
    \draw[orange] (O) circle  (3cm);
    \draw[rotate=0,->,green!50!black] (0.7,0)  arc (0:60:0.7cm);
    \draw (.8,.2) node[above,blue] {$\alpha$}; 
     \draw[green!50!black] (M)--(O)
      (M_1)--(O)--(M_2)
     (M_3)--(O)--(M_4) 
     (M)--(M_1) 
     (M_4)--(M)--(M_2); 
    \draw[dashed] (-3,-3)--(3,3)  node[above,blue] {$y=x$};
    \foreach \p/\r in {A/-45,M/45,O/-45,A'/-135,B'/-45,B/45,M_1/135,M_2/-90,M_3/-90,M_4/20} 
    \fill (\p) circle (1pt)  node[shift={(\r:3mm)},blue]{$\p$};
   \end{tikzpicture}
  } 
 } 
\end{ex}

% nguồn: Ok_dungsai_lop11_Lan1
\begin{ex}%[1D1H2-4]%[ma:8133a7dc]
 Cho góc $\alpha$ thỏa mãn  $\sin \alpha= \dfrac{1}{3}$ và $\dfrac{\pi}{2} < \alpha < \pi$. Khi đó
 \choiceTF 
  {\True $\cos \alpha =  -\dfrac{2\sqrt{2}}{3}$}
   { $\tan \alpha= \dfrac{1}{2\sqrt{2}} $}
  {\True Gọi $M$ là điểm biểu diễn  góc $\alpha$ trên đường tròn lượng giác. Khi đó
  $\mathrm{\text{sđ}}(OM,Ox)=2\pi  - \alpha+k2\pi,k \in \mathbb{Z}$}
  {\True $\sin (\alpha  + \pi)= -\dfrac{1}{3}$}
 \loigiai{ 
  \begin{itemchoice}
   \itemch \immini{Ta có $\sin \alpha= \dfrac{1}{3}$ $\Rightarrow$ ${\sin}^{2} \alpha=\dfrac{1}{9}$.\\ 
    Mà
   ${\sin}^{2} \alpha+{\cos}^{2} \alpha=1$ $\Rightarrow$ ${\cos}^{2} \alpha=1-{\sin}^{2}  \alpha=\dfrac{8}{9}$\\
   $\Leftrightarrow$ $\cos \alpha=\pm \dfrac{2\sqrt{2}}{3}$. 
   Vì $\dfrac{\pi}{2} < \alpha < \pi$ nên $\cos \alpha < 0$. Vậy $\cos \alpha=-\dfrac{2\sqrt{2}}{3}$.}{  \begin{tikzpicture}[>=stealth,line join=round,line cap=round,font=\footnotesize,scale=0.7]
   \draw [->] (-3,0)--(3,0);
   \draw[->] (0,-3)--(0,3);
  \draw (0,0) node[below right]{$O$}  coordinate (O) circle (1pt);
  \draw  (3,0) node[below]{$x$};
  \draw (0,3) node[left]{$y$}; 
  \draw[black] (0,0) circle  (2cm);
  \draw[fill=black] (160.5:2) node[left]{$M$} coordinate (M) circle  (1.2pt);
   \draw (O)--(M);
  \draw (2,0) coordinate (A)--(0:0)  coordinate (B)--(160.5:2) coordinate (C)
   pic  [draw,->]{angle=A--B--C};
  \draw (45:0.4) node[above  right]{$\alpha$};
 \end{tikzpicture}
    }
    \itemch Khi đó $\tan \alpha=\dfrac{\sin \alpha}{\cos \alpha}= -\dfrac{1}{2\sqrt{2}}$.
    \itemch
   Biểu diễn điểm $M$ trên dường tròn lượng  giác, vì $\mathrm{\text{sđ}}(Ox,OM)=\alpha$ nên
   $\mathrm{\text{sđ}}(OM,Ox)=2\pi - \alpha+k2\pi$,  $k \in \mathbb{Z}$.
   \itemch Theo tính chất của hai  góc hơn kém $\pi$ ta có $\sin (\alpha + \pi)= -\dfrac{1}{3}$.
  \end{itemchoice}
  }
\end{ex}

% nguồn: Ok_dungsai_lop11_Lan2
\begin{ex}%[1D1H2-4]%[ma:b739a5df]
	Cho  ${0< \alpha < \dfrac{\pi}{2}}$. Xét được dấu của các biểu thức sau. Khi đó
	 \choiceTF
		{\True  $\cos (\alpha+\pi) < 0$}
		{\True $\tan (\alpha-\pi) >  0$}
		{$\sin  \left(\alpha+\dfrac{2\pi}{5} \right) < 0$}
		{$\cos \left(\alpha-\dfrac{3\pi}{8}  \right) < 0$}
	 \loigiai{
		\begin{itemchoice}
			\itemch  \textbf{Đúng}.\\
			 Vì $0< \alpha < \dfrac{\pi}{2} \Rightarrow \pi < \alpha+\pi < \dfrac{3\pi}{2} \Rightarrow \cos (\alpha+\pi) < 0$.
			\itemch \textbf{Đúng}.\\ 
			Vì $0< \alpha  < \dfrac{\pi}{2} \Rightarrow-\pi < \alpha-\pi <-\dfrac{\pi}{2} \Rightarrow \tan (\alpha-\pi) > 0$.
			\itemch  \textbf{Sai}.\\
			Vì $0<  \alpha < \dfrac{\pi}{2} \Rightarrow \dfrac{2\pi}{5} < \alpha+\dfrac{2\pi}{5} < \dfrac{9\pi}{10} \Rightarrow \sin \left(\alpha+\dfrac{2\pi}{5}\right) > 0$.
			 \itemch \textbf{Sai}.\\
			Vì $0< \alpha < \dfrac{\pi}{2}  \Rightarrow-\dfrac{3\pi}{8} < \alpha-\dfrac{3\pi}{8} < \dfrac{\pi}{8} \Rightarrow \cos \left(\alpha-\dfrac{3\pi}{8}\right) > 0$.
		\end{itemchoice}
	 }
\end{ex}

% nguồn: Ok_dungsai_lop11_Lan2
\begin{ex}%[1D1H2-4]%[ma:d0054e79]
	Cho biết $\sin \alpha=\dfrac{3}{5}, \cos \alpha=-\dfrac{4}{5}$ và các biểu thức: $A=\sin \left( \dfrac{\pi}{2}-\alpha \right)+\sin (\pi +\alpha)$; $B=\cos (\pi -\alpha)+\cot \left( \dfrac{\pi}{2}-\alpha \right)$. Khi đó 
	\choiceTF 
		{\True $A=\cos \alpha  -\sin \alpha $}
		{$B=\cos  \alpha +\tan \alpha $}
		{$A+B=\dfrac{27}{20}$} 
		{\True  $A-B=-\dfrac{29}{20}$}
	\loigiai{ 
		\begin{itemchoice}
			 \itemch \textbf{Đúng.}\\
			Ta có $A=\sin \left( \dfrac{\pi}{2}-\alpha  \right)+\sin (\pi +\alpha)=\cos \alpha -\sin \alpha =-\dfrac{4}{5}-\dfrac{3}{5}=-\dfrac{7}{5}$.
			\itemch \textbf{Sai.}\\ 
			Ta có 
			\begin{align*}
				B&=\cos (\pi -\alpha)+\cot \left(  \dfrac{\pi}{2}-\alpha \right)=-\cos \alpha +\tan \alpha\\
				&=-\cos  \alpha+\dfrac{\sin \alpha}{\cos \alpha}=\dfrac{4}{5}+\dfrac{\dfrac{3}{5}}{-\dfrac{4}{5}}=\dfrac{1}{20}.
			\end{align*}
			\itemch  \textbf{Sai.}\\
			 $A+B=-\dfrac{27}{20}$.
			 \itemch \textbf{Đúng.}\\
			$A-B=-\dfrac{29}{20}$. 
		\end{itemchoice}
	 }
\end{ex}


%%%%%%%%%%%%%%%% TRẢ LỜI NGẮN %%%%%%%%%%%%%%%%

% nguồn: Ok_TraLoiNgan_Lop11_Lan2
\begin{ex}%[1D1H2-4]%[ma:49585eef]
	Biết rằng $\sin (x-\dfrac{\pi }{2})+\sin \dfrac{13\pi }{2} =\sin \left(x+\dfrac{\pi }{2} \right)$ thì giá trị đúng của $\cos x$ là 
	\shortans{0,5} 
	\loigiai{ 
		Ta có $\sin \left(x-\dfrac{\pi }{2} \right)=-\sin \left( \dfrac{\pi }{2} -x\right)=-\cos x$ và $\sin \left(x+\dfrac{\pi }{2} \right)=\cos  x$.\\
		Kết hợp với giá trị $\sin \dfrac{13\pi }{2} =\sin \left(\dfrac{\pi  }{2} +6\pi \right)=\sin \dfrac{\pi }{2} =1$.\\
		Suy  ra
		\begin{align*}
			&\sin \left(x-\dfrac{\pi }{2} \right)+\sin  \dfrac{13\pi }{2} =\sin \left(x+\dfrac{\pi }{2} \right)\\
			\Leftrightarrow&  -\cos x+1=\cos x\\
			\Leftrightarrow& \cos x=\dfrac{1}{2}. 
		\end{align*}
	} 
\end{ex}


%%%%%%%%%%%%%%%% TỰ LUẬN %%%%%%%%%%%%%%%%

% nguồn: Ok_TuanLuan_Toan11_Lan2
\begin{ex}%[1D1H2-4]%[ma:b27035c8]
	Tính giá trị của biểu thức $\sin ^210^\circ+\sin ^220^\circ+\sin  ^230^\circ+\cdots+\sin ^270^\circ+\sin ^280^\circ$.
	\loigiai{ 
		 Ta có
		\begin{eqnarray*}
			A&=&\sin ^210^\circ+\sin  ^220^\circ+\sin ^230^\circ+\cdots+\sin ^270^\circ+\sin ^280^\circ\\
			 &=&\sin ^210^\circ+\sin ^220^\circ+\sin ^230^\circ+\cdots+\cos ^220^\circ+\cos ^210^\circ \\
			&=&\sin  ^210^\circ+\cos ^210^\circ+\sin ^220^\circ+\cos ^220^\circ+\sin ^230^\circ+\cos ^230+\sin ^240^\circ+\cos ^240 \\
			 &=&4.
		\end{eqnarray*}
	 }
\end{ex}

% nguồn: Ok_TuanLuan_Toan11_Lan2
\begin{ex}%[1D1H2-4]%[ma:6f3a8013]
	 Tính giá trị của biểu thức $S=3-\sin ^290^\circ+2\cos ^260^\circ-3\tan ^245^\circ$.
	 \loigiai{
		Ta có  $S=3-\sin ^290^{^\circ}+2\cos ^260^{^\circ}-3\tan ^245^{^\circ}$ $=3-1^2+2\cdot \left(\dfrac{1}{2} \right)^2-3\cdot 1^2$ $=-\dfrac{1}{2}$.
	} 
\end{ex}

% nguồn: Ok_TuanLuan_Toan11_Lan2
\begin{ex}%[1D1H2-4]%[ma:1106e9a3]
	 Cho $2\tan \alpha-\cot \alpha=1$ với $-\dfrac{\pi}{2} < \alpha < 0$. Tính giá trị biểu thức $P=\dfrac{\tan (8\pi-\alpha)+2\cot (\pi+\alpha)}{3\tan \left(\dfrac{3\pi}{2}+\alpha\right)}$
	\loigiai{ 
		Ta có $2\tan \alpha-\cot \alpha=1\Leftrightarrow 2\tan \alpha-\dfrac{1}{\tan  \alpha}=1\Leftrightarrow \hoac{&\tan \alpha=1\\&\tan \alpha=-\dfrac{1}{2}}$.\\
		Vì  $-\dfrac{\pi}{2} < \alpha < 0$ nên $\tan \alpha < 0$.\\
		Suy ra $\tan \alpha=-\dfrac{1}{2}$ và $\cot  \alpha= \dfrac{1}{\tan \alpha}=-2$.\\
		Mà $\tan (8\pi-\alpha)=-\tan \alpha$, $\cot (\pi+\alpha)=\cot  \alpha$ và $\tan \left(\dfrac{3\pi}{2}+\alpha\right)=-\cot \alpha$.\\
		Vậy $P=\dfrac{\tan  (8\pi-\alpha)+2\cot (\pi+\alpha)}{3\tan \left(\dfrac{3\pi}{2}+\alpha\right)}=\dfrac{-\tan \alpha+2 \cot \alpha}{-3 \cot \alpha}=\dfrac{\dfrac{1}{2}-4}{6}=\dfrac{-7}{12}.$
	} 
\end{ex}

% nguồn: Ok_TuanLuan_Toan11_Lan2
\begin{ex}%[1D1H2-4]%[ma:265f718e]
	Tính giá trị  của biểu thức
	\begin{align*}
		M=&\cos ^210^\circ+\cos ^220^\circ+\cos ^230^\circ+\cdots+\cos ^2170^\circ+\cos  ^2180^\circ.
	\end{align*}
	 \loigiai{
		Áp dụng công thức $\cos \alpha=\cos \left(180^\circ-\alpha \right)$ và $\cos ^2\alpha+\sin ^2\alpha=1$  ta có
		\begin{eqnarray*}
			M&=&\cos ^210^\circ+\cos ^220^\circ+\cos  ^230^\circ+\cdots+\cos ^2170^\circ+\cos ^2180^\circ\\
			&=&\cos ^210^\circ+\cos ^220^\circ+\cdots+\cos ^280^\circ+\cos ^290^\circ+\cos ^280^\circ+\cdots++\cos ^220^\circ+\cos  ^210^\circ+\cos ^290^\circ\\
			&=&2\left(\cos ^210^\circ+\cos ^220^\circ+\cos ^230^\circ+\cdots+\cos ^280^\circ+\cos  ^290^\circ\right)\\
			 &=&2\left(\sin ^280^\circ+\cdots+\sin ^250^\circ+\cos ^250^\circ+\cdots+\cos ^280^\circ+\cos ^290^\circ\right)=8.
		\end{eqnarray*}
	} 
\end{ex}

% nguồn: Ok_TuanLuan_Toan11_Lan2
\begin{ex}%[1D1H2-4]%[ma:ae443ed8]
	Bằng cách sử dụng giá trị lượng giác của  các góc liên quan đặc biệt, hãy tính
	\begin{enumerate}
		\item  $\cos \left(\dfrac{-15\pi}{4}\right)$;
		\item $\cot  \left(-675^\circ \right)$;
		\item $\sin  \dfrac{11\pi}{3}$.
	\end{enumerate}
	 \loigiai{
		\begin{enumerate}
			\item $\cos \left(\dfrac{-15\pi}{4}\right) =\cos \dfrac{15\pi}{4} =\cos \left(3\pi +\dfrac{3\pi}{4} \right)=-\cos \dfrac{3\pi}{4} =\cos \left(\pi -\dfrac{3\pi}{4} \right)=\cos \dfrac{\pi}{4} =\dfrac{\sqrt{2}}{2}  $.
			\item $\cot (-675^\circ)=\cot (45^\circ -720^\circ)=\cot 45^\circ  =1$.
			\item $\sin  \dfrac{11\pi}{3} =\sin \left(3\pi +\dfrac{2\pi}{3} \right)=-\sin \dfrac{2\pi}{3} =-\sin \left(\pi -\dfrac{2\pi}{3} \right)=-\sin \dfrac{\pi}{3} =-\dfrac{\sqrt{3}}{2} $.
		\end{enumerate}
	} 
\end{ex}

% nguồn: Ok_TuanLuan_Toan11_Lan2
\begin{ex}%[1D1H2-4]%[ma:dcdb8d6f]
	Rút gọn biểu thức  $D=\sin \left(\dfrac{5\pi}{2}-\alpha \right)+\cos \left(13\pi+\alpha \right)-3\sin \left(\alpha-5\pi \right)$.
	\loigiai{ 
		Ta  có
		\begin{align*}
			D&=\sin \left(\dfrac{5\pi}{2}-\alpha \right)+\cos  \left(13\pi+\alpha \right)-3\sin \left(\alpha-5\pi \right)\\
			&=\sin \left(\dfrac{\pi}{2}-\alpha \right)+\cos \left(\pi+\alpha  \right)+3\sin \left(\pi-\alpha \right)=\cos \alpha-\cos \alpha+3\sin \alpha=3\sin \alpha.
		\end{align*}
	} 
\end{ex}

% nguồn: Ok_TuanLuan_Toan11_Lan2
\begin{ex}%[1D1H2-4]%[ma:7f70476f]
	Hãy  thực hiện các yêu cầu trong mỗi trường hợp sau
	\begin{enumerate}
		\item Biểu diễn $\sin \left(-\dfrac{29\pi}{3} \right)$ qua giá trị lượng giác của góc có số đo từ 0 đến $\dfrac{\pi}{4} $; 
		\item Biểu diễn tan $973^\circ $ qua giá trị  lượng giác của góc có số đo từ $0^\circ $ đến $45^\circ $.
	\end{enumerate}
	\loigiai{ 
		\begin{enumerate}
			\item  Ta có
			\begin{eqnarray*}
				\sin \left(-\dfrac{29\pi}{3} \right) &=&-\sin \left(\dfrac{29\pi}{3} \right)=-\sin  \left(\dfrac{5\pi}{3} +8\pi \right)\\
				&=& -\sin \left(\dfrac{2\pi}{3} +\pi  \right)=\sin \dfrac{2\pi}{3}\\
				&=&\cos \left(\dfrac{\pi}{2}  -\dfrac{2\pi}{3} \right)=\cos \left(-\dfrac{\pi}{6} \right)=\cos \dfrac{\pi}{6} .
			\end{eqnarray*}
			\item $\tan  973^\circ =\tan (73^\circ +5\cdot180^\circ)=\tan 73^\circ =\tan (90^\circ -17^\circ)=\cot 17^\circ $.
	\end{enumerate}}
\end{ex}
